%!TEX root = ../Master_template.tex
\chapter{Experiments}

All experiments were executed on the university HPC cluster (STUD partition). Each experiment
was allocated 4 CPU cores and a single NVIDIA GeForce RTX 2070 GPU with 8\,GB of memory.

\section{Experimental Setup}
\label{sec:exp_setup}

\textbf{Problem and Dataset.} All methods are evaluated using a real-world dataset from Beijing, which was also used by the
authors of the baseline paper \citep{zhang2023}. The dataset consists of 100 locations
(including the depot), and travel times are recorded for every ordered pair $(i, j)$, yielding
10{,}000 city pairs in total. The 24-hour day is divided into 2-hour intervals (giving 12 time
bins per arc), with each interval storing a normalised travel time. Rather than using raw
(actual) values, the travel times are preprocessed using cubic spline interpolation. Full
details of this procedure are described in Section~\ref{subsec:data_gen}. Out of these 100
locations, this study focuses on four customer sets ($C$):

\begin{itemize}
    \item $C = 19$: used for comparisons in the experiments.
    \item $C = 20$: used for hyperparameter optimisation.
    \item $C = 29$: intermediate scalability check and graph dependency experiments.
    \item $C = 49$: used for larger-scale and scalability experiments.
\end{itemize}

Each training instance uses a random subset of the 100 cities of size graph\_size $+ 1$
(including the depot), selected by random scoring and sorting as described in
Section~\ref{subsec:data_gen}. Nodes are represented as one-hot encoded vectors of dimension
100 rather than raw coordinates, since the model operates on time-dependent travel times
rather than Euclidean distances.

\begin{figure}[H]
    \centering
    \includegraphics[width=0.75\textwidth]{figs/location_map_of_benchmark_dataset.png}
    \caption{Location map of the benchmark dataset (Beijing). The depot is located at the
    bottom-right of the map. Reproduced from \citep{zhang2023}.}
    \label{fig:beijing_map}
\end{figure}

\textbf{Models and variants.} As discussed in Section~\ref{sec:extensions}, this chapter evaluates:

\begin{itemize}
    \item Four different step-aware decoder configurations with different combinations of
    features.
    \item Three different heuristic model blending configurations using cost-aware gating.
    \item Three different time-sliced window encoding and safe refreshing mechanism
    configurations.
\end{itemize}

\textbf{Heuristic Baselines.} The proposed model's performance is compared with the following classical heuristics:
Nearest Neighbour (NN), Greedy Edge, Randomised Nearest Neighbour (NNR),
look-ahead-1 (NN+1), and 2-opt initialised with Nearest Neighbour.

For a fair comparison, all heuristics are evaluated on the same problem instances and use the
same cubic spline interpolation method.

\textbf{Training.} All models are trained with policy gradient reinforcement learning with a rollout baseline
using the following hyperparameters:

\begin{table}[H]
\centering
\caption{Hyperparameters explored and selected values.}
\label{tab:hyperparams_explored}
\begin{adjustbox}{max width=\textwidth}
\begin{tabular}{lll}
\toprule
\textbf{Hyperparameter} & \textbf{Values Explored} & \textbf{Selected Value} \\
\midrule
Batch size          & 32, 256, 512                       & 256 (see Section~\ref{sec:batch_size_stability}) \\
Epochs              & 50, 100                            & 50 ($C = 19/20$), 100 ($C = 29/49$) \\
Learning rate       & 0.00005, 0.0001, 0.00015, 0.0002  & 0.0001 (see Section~\ref{sec:hyperparam_opt}) \\
LR decay            & 0.99, 0.995, 0.998                 & 0.995 (see Section~\ref{sec:hyperparam_opt}) \\
Step-MLP hidden dim & 32, 64, 128                        & 64 (see Section~\ref{sec:hyperparam_opt}) \\
Random seeds        & 1234, 5678, 9012                   & All three (averaged) \\
\bottomrule
\end{tabular}
\end{adjustbox}
\end{table}

To ensure robust results, three predefined random seeds are used, and all reported numbers are
averages of the three random seeds. Sampling is used for training, and greedy decoding is used
for both the baseline model and evaluation, unless otherwise specified. The user may optionally
select beam search for evaluation, and those results are reported separately in
Section~\ref{sec:beam_search}.

\section{Baseline Model (M1) Performance}
\label{sec:baseline_performance}

Before evaluating the proposed modifications, the performance of the M1 baseline is first
established as an unmodified reference point. This baseline uses the full-day traffic profile,
without decoder MLP variants or cost-aware gating. All results use greedy decoding and are
reported as mean test cost in minutes.

\textbf{Effect of Training Duration ($C = 19$).} Using the architecture described in Chapter~3, the baseline model (M1) is trained for
$C = 19$ nodes with 50 and 100 epochs. Table~\ref{tab:baseline_combined} shows that training
for 100 epochs improves performance by approximately 9.2 minutes (574.73 vs 565.50 min),
with significantly reduced variance (10.78 vs 0.58). Based on this finding, all experiments
on larger graph sizes ($C = 29$, $C = 49$) use 100 training epochs as the default, while
experiments focused on feature selection and hyperparameter search at $C = 19$ and
$C = 20$ use 50 epochs to reduce computational cost and obtain results more quickly.

\begin{table}[H]
\centering
\caption{Baseline M1 performance across training epochs and graph sizes.}
\label{tab:baseline_combined}
\begin{adjustbox}{max width=\textwidth}
\begin{tabular}{lll}
\toprule
\textbf{Graph Size} & \textbf{Epochs} & \textbf{Mean $\pm$ SD (min)} \\
\midrule
$C = 19$ & 50  & 574.73 $\pm$ 10.78 \\
$C = 19$ & 100 & 565.50 $\pm$ 0.58 \\
$C = 29$ & 100 & 766.79 $\pm$ 7.39 \\
$C = 49$ & 100 & 1095.93 $\pm$ 73.34 \\
\bottomrule
\end{tabular}
\end{adjustbox}
\end{table}

These baseline results serve as the comparison target throughout subsequent sections.

\section{Step-MLP Variants}
\label{sec:step_mlp_variants}

This section evaluates the proposed Step-MLP and related decoder enhancements. All
experiments use $C = 19$ with 50 epochs unless otherwise stated, with the reference baseline
of $574.73 \pm 10.78$ minutes (Table~\ref{tab:baseline_combined}). The Step-MLP adds a
context nudge to the decoder using step-aware features passed through a three-layer MLP
(Section~\ref{subsec:step_mlp}).

\textbf{Feature Selection for Step-MLP.} Individual features (step ratio, linear time, depot distance, tour length, mean dist
unvisited, last 3 visited nodes) were first evaluated independently. No single feature
outperformed the baseline, ranging from 577.9 to 581.5 minutes, compared to 574.7 minutes
for the baseline, though linear time and mean distance to unvisited nodes performed best
among them.

Several predefined feature combinations were then evaluated (Table~\ref{tab:feature_presets_exp}).
The v2\_light preset with linear time encoding achieves the lowest mean cost (573.80 min).
However, this difference of 0.93 minutes lies well within one standard deviation of the baseline
(574.73 $\pm$ 10.78 min) and cannot be considered a statistically significant improvement at
this scale. Notably, v2\_light exhibits substantially reduced variance ($\pm$7.22 vs.\ $\pm$10.78),
indicating more stable training across seeds. Higher-dimensional feature
sets (129 to 132 dimensions) do not offer consistent improvements, suggesting embedding-based
features introduce noise at this problem size.

\begin{table}[H]
\centering
\caption{Feature preset comparison ($C = 19$, 50 epochs).}
\label{tab:feature_presets_exp}
\begin{adjustbox}{max width=\textwidth}
\begin{tabular}{llllll}
\toprule
\textbf{Feature Set} & \textbf{Baseline} & \textbf{v1 light (4)}
    & \textbf{minimal (129)} & \textbf{v1 (132)} & \textbf{v2 light (linear, 5)} \\
\midrule
Mean $\pm$ SD (min) & 574.73 $\pm$ 10.78 & 580.00 $\pm$ 3.42
    & 575.60 $\pm$ 8.00 & 574.95 $\pm$ 7.36 & \textbf{573.80 $\pm$ 7.22} \\
\bottomrule
\end{tabular}
\end{adjustbox}
\end{table}

A leave-one-out ablation on v2\_light, which has five features, (Table~\ref{tab:ablation_v2light})
reveals that linear time is the most critical feature, as removing it increases cost by 7.5
minutes. Replacing linear time with sinusoidal encoding also degrades performance (578.2 vs.\
573.8 min), likely because travel time bins represent a single non-repeating operational day.
Based on its stable training behaviour and lowest mean cost among all tested presets, v2\_light
with linear time encoding is selected for all subsequent experiments.

\begin{table}[H]
\centering
\caption{Leave-one-out ablation of v2\_light ($C = 19$, 50 epochs).}
\label{tab:ablation_v2light}
\begin{adjustbox}{max width=\textwidth}
\begin{tabular}{lll}
\toprule
\textbf{Removed Feature} & \textbf{Mean Cost (min)} & \textbf{$\Delta$ vs.\ Full v2\_light} \\
\midrule
None (full v2\_light)  & 573.80 & - \\
linear\_time           & 581.34 & $+7.49$ \\
mean\_dist\_unvisited  & 578.80 & $+4.95$ \\
step\_ratio            & 577.96 & $+4.11$ \\
tour\_length           & 577.12 & $+3.27$ \\
depot\_distance        & 574.55 & $+0.70$ \\
\bottomrule
\end{tabular}
\end{adjustbox}
\end{table}

\textbf{Temperature MLP and Decoder MLP Variants.} Beyond the Step-MLP context nudge, three additional mechanisms were evaluated: Temperature
MLP (scales attention logits), Option C (pre-attention MLP), and Option A (post-attention
MLP). Table~\ref{tab:decoder_mlp_c19} summarises results at $C = 19$.

\begin{table}[H]
\centering
\caption{Decoder MLP variant results ($C = 19$, 50 epochs).}
\label{tab:decoder_mlp_c19}
\begin{adjustbox}{max width=\textwidth}
\begin{tabular}{ll}
\toprule
\textbf{Configuration} & \textbf{Mean $\pm$ SD (min)} \\
\midrule
Baseline             & 574.73 $\pm$ 10.78 \\
Temp-MLP standalone  & \textbf{568.68 $\pm$ 2.12} \\
Option C only        & 570.39 $\pm$ 1.85 \\
Option A only        & 570.32 $\pm$ 2.10 \\
Step-MLP + Temp-MLP  & 576.00 $\pm$ 5.59 \\
Option C + Temp-MLP  & 570.05 $\pm$ 0.98 \\
\bottomrule
\end{tabular}
\end{adjustbox}
\end{table}

At $C = 19$, Temp-MLP standalone achieves the best performance (568.68 min), improving by
approximately 6 minutes with notably reduced variance. Option A and Option C individually
outperform the baseline by almost 4 minutes. However, combining components generally does
not yield additional gains.

\textbf{Scaling to $C = 49$.} Table~\ref{tab:decoder_mlp_c49} presents results at the
larger graph size with 100 epochs.

\begin{table}[H]
\centering
\caption{Decoder MLP variant comparison ($C = 49$, 100 epochs, batch size = 256).}
\label{tab:decoder_mlp_c49}
\begin{adjustbox}{max width=\textwidth}
\begin{tabular}{lll}
\toprule
\textbf{Configuration} & \textbf{Mean $\pm$ SD (min)} & \textbf{$\Delta$ vs.\ Baseline} \\
\midrule
Baseline                                 & 1095.93 $\pm$ 73.34 & -- \\
Temp-MLP only                            & \textbf{1041.57 $\pm$ 19.61} & $-4.96\%$ \\
Step-MLP + Temp-MLP                      & 1049.43 $\pm$ 22.91 & $-4.24\%$ \\
Option C only                            & 1061.66 $\pm$ 37.07 & $-3.13\%$ \\
Step-MLP + Option C + Temp-MLP$^\dagger$ & 1062.48 $\pm$ 25.23 & $-3.05\%$ \\
Step-MLP (v2\_light)                     & 1068.77 $\pm$ 22.56 & $-2.48\%$ \\
Option A only                            & 1158.85 $\pm$ 67.32 & $+5.74\%$ \\
\bottomrule
\end{tabular}
\end{adjustbox}
\end{table}

\noindent{\footnotesize $^\dagger$Evaluation batch size reduced to 256 (default: 1024) due to
CUDA memory constraints.}

\medskip

At $C = 49$, Temp-MLP alone yields the strongest performance (1041.57 min), delivering a
54-minute (4.96\%) improvement over the baseline. Option A performs worse than the baseline,
suggesting that post-attention glimpse refinement is not beneficial and that it is better to
retain the linear transformation rather than a non-linear one. Combining components does not
provide meaningful gains over the best standalone variant.

\begin{figure}[H]
\centering
\includegraphics[width=0.95\textwidth]{figs/val_cost_curves_temp_baseline.png}
\caption{Mean validation cost $\pm$ standard deviation across three seeds (1234, 5678, 9012)
over 300 epochs at $C = 49$ for Temp-MLP (solid blue) and Baseline (dashed orange). The
shaded regions represent $\pm$ one standard deviation. The Temp-MLP mean remains consistently
below the baseline throughout training, with a stable gap from approximately epoch 30 onwards.}
\label{fig:val_cost_300ep}
\end{figure}

Due to computational and time limitations, the experiments in the main study were restricted
to 100 epochs. However, it is still worth investigating the impact of longer training.
Therefore, both the baseline and Temp-MLP were trained for 300 epochs across all three seeds,
which corresponds to triple the default compute budget used in the main experiments.
Figure~\ref{fig:val_cost_300ep} shows the mean validation cost across the three seeds over
the full 300 epochs.

The mean curve confirms that Temp-MLP maintains a consistent advantage over the baseline
throughout the training process. Around epochs 20--30, the gap begins to open and remains
almost stable for the rest of training. Importantly, Temp-MLP shows comparatively lower
variance across seeds than the baseline, indicating more consistent training behaviour across
different random initialisations.

Neither model is fully converged at epoch 300, as both curves are still slowly declining.
However, the rate of improvement clearly diminishes in the later epochs, suggesting that
additional training would likely provide only marginal gains. The relative ordering between
the two models is also unlikely to change.

\textbf{Summary and Recommendation.}

\begin{figure}[H]
\centering
\includegraphics[width=0.95\textwidth]{figs/figure_4_5_decoder_mlp_final.png}
\caption{Summary of decoder MLP variant performance at $C = 19$ and $C = 49$.}
\label{fig:decoder_mlp_summary}
\end{figure}

The Temp-MLP produces the most significant and consistent improvement at both graph sizes,
with notably low variance. At $C = 49$, it achieves a 4.96\% cost reduction. Option A
(post-attention MLP) degrades performance at larger scales, suggesting that post-attention
glimpse refinement is not beneficial and that it is better to retain the linear transformation
rather than a non-linear one. Since the individual components outperform their combinations, Temp-MLP standalone is chosen
as the best configuration for the final model comparison (Section~\ref{sec:heuristic_comparison}).
As discussed earlier, due to computational and time limitations, the experiments in the main
study were restricted to 100 epochs, where a single 100-epoch run at $C = 49$ takes
approximately 360--650 minutes per seed. While 300-epoch runs confirmed that the Temp-MLP
advantage is stable and persistent, running all configurations for 300 epochs was not feasible
within the available compute budget. The 100-epoch results are therefore used consistently
across all experiments to ensure fair and comparable evaluation. As shown in
Figure~\ref{fig:val_cost_300ep}, the relative ordering between Temp-MLP and the baseline is
stable from epoch 30 onwards, providing confidence that the conclusions drawn from 100-epoch
experiments are reliable despite the limited training budget.

\section{Batch Size and Training Stability}
\label{sec:batch_size_stability}

During initial experiments at $C = 49$, CUDA memory errors were encountered, and as a
solution the batch size was reduced (original size was 512). This reduction introduced training instability, so
batch sizes of 32 and 256 were compared (Table~\ref{tab:batch_size}).

\begin{table}[H]
\centering
\caption{Batch size effect on test cost and training stability ($C = 49$, 100 epochs).}
\label{tab:batch_size}
\begin{adjustbox}{max width=\textwidth}
\begin{tabular}{llll}
\toprule
\textbf{Model} & \textbf{Batch} & \textbf{Mean $\pm$ SD (min)} & \textbf{Best Epoch} \\
\midrule
Baseline           & 32  & 1086.67 $\pm$ 13.84 & 9 \\
Baseline           & 256 & 1095.93 $\pm$ 73.34 & 86 \\
Step-MLP v2\_light & 32  & 1072.74 $\pm$ 43.31 & 35 \\
Step-MLP v2\_light & 256 & 1068.77 $\pm$ 22.56 & 97 \\
\bottomrule
\end{tabular}
\end{adjustbox}
\end{table}

Batch size 32 produces early best epochs (epochs 9 and 35), indicating training instability.
In contrast, batch size 256 converges more stably, with the best epoch found at epochs 86
to 97. Therefore, \textbf{batch size 256} is used for all remaining experiments.

\section{Hyperparameter Optimisation}
\label{sec:hyperparam_opt}

After selecting the best-performing Step-MLP feature configuration (v2\_light with linear
time, Section~\ref{sec:step_mlp_variants}), the next step is to identify the best
hyperparameter combination for training. All experiments in this section use the Step-MLP
v2\_light configuration, $C = 20$, batch size 256, and 50 epochs.

Three hyperparameters are tuned sequentially instead of computationally expensive grid
search: Step-MLP hidden dimension, learning rate, and learning rate decay.
Figure~\ref{fig:hyperparams} shows the results of the hyperparameter sweep.

For hidden dimension, 64 yields the best performance (591.99 minutes) compared to 32
(600.49 minutes) and 128 (603.06 minutes).

Using the selected hidden dimension of 64, learning rate 0.0001 achieves the best result
(591.99 minutes), while other values result in performance degradation of 4 to 12 minutes.

The learning rate decay is applied multiplicatively to the learning rate at the end of each
epoch. A decay of 0.995 provides a slight (approximately one minute) advantage compared to
the alternatives.

The finalised hyperparameter combination is step\_mlp\_dim\,=\,64, learning\_rate\,=\,0.0001,
and lr\_decay\,=\,0.995. Table~\ref{tab:hyperparam_results} summarises the selected values.

\begin{table}[H]
\centering
\caption{Selected hyperparameters ($C = 20$, Step-MLP v2\_light, batch size 256).}
\label{tab:hyperparam_results}
\begin{adjustbox}{max width=\textwidth}
\begin{tabular}{llll}
\toprule
\textbf{Hyperparameter} & \textbf{Values Tested} & \textbf{Selected}
    & \textbf{Mean $\pm$ SD (min)} \\
\midrule
Step-MLP dim  & 32, 64, 128               & 64    & 591.99 $\pm$ 3.02 \\
Learning rate & 5e-5, 1e-4, 1.5e-4, 2e-4 & 1e-4  & 591.99 $\pm$ 3.02 \\
LR decay      & 0.99, 0.995, 0.998        & 0.995 & 591.99 $\pm$ 3.02 \\
\bottomrule
\end{tabular}
\end{adjustbox}
\end{table}

\begin{figure}[H]
\centering
\includegraphics[width=0.95\textwidth]{figs/figure_4_6_hyperparams.png}
\caption{Hyperparameter sweep results: Step-MLP hidden dimension (left), learning rate
(centre), and learning rate decay (right). All experiments use $C = 20$, batch size 256,
50 epochs.}
\label{fig:hyperparams}
\end{figure}

\section{Cost-Aware Gating}
\label{sec:cost_aware_exp}

This section investigates the effect of incorporating cost-aware gating into the M1 model,
which adds a heuristic bias to the decoder's logits (Section~\ref{subsec:cost_aware}). Three
heuristic functions are evaluated: Nearest Neighbour (NN), Nearest Neighbour Random (NNR),
and one-step lookahead (NN+1). The reference baseline for $C = 19$ is 565.50 minutes and
for $C = 49$ is 1095.93 minutes (Table~\ref{tab:baseline_combined}).

\textbf{Fixed Lambda ($C = 19$).} Figure~\ref{fig:lambda_sensitivity} shows the effect of various $\lambda$ values for the NN
heuristic at $C = 19$. The optimal range is $\lambda \in [0.25, 0.75]$, and all values
within this range outperform the baseline. When $\lambda$ exceeds 1.0, performance degrades
substantially, indicating that over-weighting the heuristic suppresses the learned attention
signal.

\begin{figure}[H]
\centering
\includegraphics[width=0.80\textwidth]{figs/figure_4_7_lambda_sensitivity.png}
\caption{Fixed lambda sensitivity with NN heuristic ($C = 19$, seed 1234).}
\label{fig:lambda_sensitivity}
\end{figure}

Table~\ref{tab:fixed_lambda_c19} compares the three heuristic functions with the previously
identified fixed $\lambda$ value. Ideally, the best fixed $\lambda$ value would be identified
separately for each heuristic. However, due to computational limitations, the NN best fixed
$\lambda$ value is applied to the other two methods as well. The NN heuristic achieves the
best performance among the three, yet is unable to beat the baseline at this graph size.

\begin{table}[H]
\centering
\caption{Fixed lambda results by heuristic ($C = 19$, 100 epochs).}
\label{tab:fixed_lambda_c19}
\begin{adjustbox}{max width=\textwidth}
\begin{tabular}{lllll}
\toprule
\textbf{Metric} & \textbf{Baseline} & \textbf{NN} & \textbf{NNR} & \textbf{NN+1} \\
\midrule
Mean $\pm$ SD (min) & 565.50 $\pm$ 0.58 & 566.80 $\pm$ 0.58 & 572.90 $\pm$ 3.70 & 570.21 $\pm$ 3.24 \\
\bottomrule
\end{tabular}
\end{adjustbox}
\end{table}

Learnable $\lambda$ was also tested using gradient descent but underperformed fixed $\lambda$
across all heuristics (575.43 min for NN vs.\ 566.80 min with fixed $\lambda$). Nonlinear
transformations (exponential and piecewise) were also tested but failed to outperform the
linear approach (575.43 min and 567.83 min, respectively, vs.\ baseline 565.50 min).

\textbf{Scaling to $C = 49$.} Since the cost-aware gating mechanism failed to beat the baseline at $C = 19$, it is worth
investigating how it performs at $C = 49$. The best-performing configuration at $C = 19$
(NN heuristic, fixed $\lambda = 0.25$) was used to evaluate $C = 49$ (Table~\ref{tab:cost_aware_c49}).

\begin{table}[H]
\centering
\caption{Cost-aware gating results ($C = 49$, 100 epochs, batch size 256, fixed $\lambda =
0.25$).}
\label{tab:cost_aware_c49}
\begin{adjustbox}{max width=\textwidth}
\begin{tabular}{lll}
\toprule
\textbf{Configuration} & \textbf{Mean $\pm$ SD (min)} & \textbf{$\Delta$ vs.\ Baseline} \\
\midrule
Baseline                    & 1095.93 $\pm$ 73.34 & --- \\
Cost-aware only (NN)        & \textbf{1055.75 $\pm$ 25.14} & $-3.67\%$ \\
Option C + Cost-aware (NN)  & 1068.17 $\pm$ 26.05 & $-2.53\%$ \\
Temp-MLP + Cost-aware (NN)  & 1062.78 $\pm$ 35.90 & $-3.02\%$ \\
\bottomrule
\end{tabular}
\end{adjustbox}
\end{table}

Unlike at $C = 19$, cost-aware gating leads to a significant improvement at $C = 49$.
Cost-aware gating with the NN heuristic alone produces an improvement of approximately 40
minutes (3.67\%) compared to the baseline. Combining with other decoder components does not
provide additional gains, so standalone cost-aware gating yields the best performance. This
suggests that cost-aware gating is particularly beneficial at larger problem sizes, as the
model struggles with larger problems and cost-aware gating provides additional guidance
through cost information to produce better results.

\textbf{Summary and Recommendation.} At $C = 19$, cost-aware gating does not outperform the baseline. Learnable $\lambda$
underperformed fixed $\lambda$, and nonlinear transformations failed to improve upon the
linear approach.

At $C = 49$, cost-aware gating produces a notable improvement (3.67\%) compared to the
baseline. Combining with other decoder components does not provide additional gains.

The best cost-aware configuration is: fixed $\lambda = 0.25$, NN heuristic, linear
transformation.

\section{Time Slicing with Safe Refresh}
\label{sec:time_slicing_exp}

The baseline M1 encoder feeds the full 24-hour traffic data. Time slicing replaces this with
a forward window of $W$ time bins starting from a given start bin, thereby reducing the input
dimensionality. When the decoder's current time approaches the end of the window, a safe
refresh mechanism re-encodes (Section~\ref{subsec:time_slicing}).

Before presenting the experimental results, the motivating hypotheses are introduced. These
hypotheses were formulated based on the structure of the TDTSP and the architecture of the
baseline model (Section~\ref{subsec:time_slicing}):

\begin{itemize}
    \item \textbf{H1:} Shorter time windows improve solution quality by reducing noise.
    \item \textbf{H2:} More frequent encoder refresh improves performance.
    \item \textbf{H3:} Optimal configurations depend on graph size.
    \item \textbf{H4:} Start time has a meaningful impact on solution quality.
    \item \textbf{H5:} Time slicing reduces computational cost with small accuracy loss.
\end{itemize}

These hypotheses are assessed in the summary and recommendations.

The experiments evaluate three configuration choices: window size ($W$), refresh strategy,
and start time. Results are presented across graph sizes $C = 19$, 20, 29, and 49. The
baseline references are 574.73 minutes ($C = 19$, 50 epochs), 766.79 minutes ($C = 29$),
and 1095.93 minutes ($C = 49$).

\textbf{Window Size Comparison.} Four window sizes ($W = 3, 6, 9$, and forward window) were evaluated across three graph
sizes with fixed start bin $= 3$ and one-time refresh. Figure~\ref{fig:window_size}
summarises the results.

\begin{figure}[H]
\centering
\includegraphics[width=0.95\textwidth]{figs/figure_4_9_window_size.png}
\caption{Mean test cost vs.\ window size $W$ for $C = 19$, $C = 20$, and $C = 29$ (start
bin = 3, one-time refresh, 50 epochs).}
\label{fig:window_size}
\end{figure}

At $C = 19$, no window improves upon the baseline, though $W = 9$ performs best among the
windows tested. At $C = 20$, almost all time slicing configurations achieve performance at or
better than the baseline. Among them, $W = 3$ achieves the best mean cost (590.50 min) with
the lowest variance. To continue investigating the effect of time windows at larger problem
sizes, consider $C = 29$. At $C = 29$, $W = 9$ yields better performance than the baseline,
as narrower windows (3 and 6) appear insufficient for larger tours. The optimal window size
thus depends on graph size.

\textbf{Refresh Strategy Comparison.} Four refresh strategies were evaluated: one-time, periodic (intervals 0.25 and 0.50),
k-moves lookahead buffer, and combined. Table~\ref{tab:refresh_summary} summarises the
best-performing strategy for each graph size.

\begin{table}[H]
\centering
\caption{Best refresh strategy per graph size.}
\label{tab:refresh_summary}
\begin{adjustbox}{max width=\textwidth}
\begin{tabular}{lllll}
\toprule
\textbf{Graph Size} & \textbf{Window} & \textbf{Best Strategy}
    & \textbf{Mean $\pm$ SD (min)} & \textbf{$\Delta$ vs.\ Baseline} \\
\midrule
$C = 19$ & $W = 9$ & Periodic (0.25) & 572.55 $\pm$ 2.86 & $-0.38\%$ \\
$C = 20$ & $W = 3$ & One-time        & 590.50 $\pm$ 3.08 & $-0.91\%$ \\
$C = 29$ & $W = 9$ & One-time        & 763.48 $\pm$ 1.11 & $-0.43\%$ \\
\bottomrule
\end{tabular}
\end{adjustbox}
\end{table}

All of the best refreshing strategies for graph sizes 19, 20, and 29 beat the baseline. At
$C = 19$ and 20, periodic refresh (interval = 0.25) performs best. At $C = 29$, periodic
refresh (0.50) just beats the baseline, but one-time performed best. One-time refresh is
selected for future analysis due to its simplicity and competitive performance with low
variance.

\textbf{Start Time Sensitivity.} Start bins 0, 3, 6, and 9 were evaluated for each graph size. Start bin 3 consistently
performs best across all graph sizes. Start bin 0 provides insufficient context, while bins
6 and 9 may introduce irrelevant next-day data when combined with larger windows.

\textbf{Validation at $C = 49$.} Using the best configuration ($W = 3$, start bin = 3), time slicing is validated at
$C = 49$ with 100 epochs (Table~\ref{tab:time_slicing_c49}).

\begin{table}[H]
\centering
\caption{Time slicing validation ($C = 49$, $W = 3$, start bin = 3, 100 epochs).}
\label{tab:time_slicing_c49}
\begin{adjustbox}{max width=\textwidth}
\begin{tabular}{lll}
\toprule
\textbf{Strategy} & \textbf{Mean $\pm$ SD (min)} & \textbf{$\Delta$ vs.\ Baseline} \\
\midrule
Baseline        & 1095.93 $\pm$ 73.34 & -- \\
One-time        & 1094.39 $\pm$ 43.17 & $-0.14\%$ \\
Periodic (0.25) & \textbf{1082.81 $\pm$ 26.19} & $-1.20\%$ \\
\bottomrule
\end{tabular}
\end{adjustbox}
\end{table}

At $C = 49$, periodic (0.25) achieves the best performance of 1082.81 minutes, an
improvement of 13 minutes (1.20\%) over the baseline. More frequent encoder updates
contribute to better solution quality at larger graph sizes.

\textbf{Computational Efficiency.} Table~\ref{tab:efficiency} compares training time and solution quality.

\begin{table}[H]
\centering
\caption{Computational efficiency: time slicing vs.\ baseline.}
\label{tab:efficiency}
\begin{adjustbox}{max width=\textwidth}
\begin{tabular}{llll}
\toprule
\textbf{Graph Size} & \textbf{Best Config} & \textbf{$\Delta$ Cost}
    & \textbf{Time Overhead} \\
\midrule
$C = 19$ & $W=9$, Periodic (0.25) & $-0.38\%$ & $+2.49\%$ \\
$C = 20$ & $W=3$, one-time        & $-0.91\%$ & $+1.70\%$ \\
$C = 29$ & $W=9$, one-time        & $-0.43\%$ & $+8.33\%$ \\
$C = 49$ & $W=3$, Periodic (0.25) & $-1.19\%$ & $+10.72\%$ \\
\bottomrule
\end{tabular}
\end{adjustbox}
\end{table}

Time slicing improves solution quality through regular encoder refreshing and the exclusion
of noisy, irrelevant data. However, it adds training overhead due to the refresh mechanism
checking window boundaries and re-running the encoder. Rather than saving computation, time
slicing trades additional training time for better solutions.

\textbf{Summary and Recommendation.} Table~\ref{tab:time_slicing_summary} consolidates the best time-slicing configuration and
its outcome at each graph size. The five hypotheses defined at the outset are assessed below.

\begin{table}[H]
\centering
\caption{Time slicing summary - best configuration per graph size.}
\label{tab:time_slicing_summary}
\begin{adjustbox}{max width=\textwidth}
\begin{tabular}{lllll}
\toprule
\textbf{Graph Size} & \textbf{Best Config} & \textbf{Mean Cost (min)}
    & \textbf{$\Delta$ vs.\ Baseline} & \textbf{Time Overhead} \\
\midrule
$C = 19$ & $W=9$, start=3, Periodic (0.25) & 572.55  & $-0.38\%$ & $+2.49\%$ \\
$C = 20$ & $W=3$, start=3, one-time        & 590.50  & $-0.91\%$ & $+1.70\%$ \\
$C = 29$ & $W=9$, start=3, one-time        & 763.48  & $-0.43\%$ & $+3.8\%$ \\
$C = 49$ & $W=3$, start=3, Periodic (0.25) & 1082.81 & $-1.19\%$ & $+10.72\%$ \\
\bottomrule
\end{tabular}
\end{adjustbox}
\end{table}

The five hypotheses defined at the outset of this section are now assessed in light of the
experimental findings.

\textbf{H1 (Shorter windows improve quality) - Mostly rejected.} The forward window
($W = 9$) outperforms narrow windows in 2/3 graph sizes ($C = 19$ and $C = 29$). This
suggests that feeding the encoder with current and all remaining future data is more
beneficial.

\textbf{H2 (Frequent refresh improves performance) - Mostly confirmed.} One-time refresh
was selected as the practical choice because of its simplicity, but periodic refresh wins
in 3/4 cases (except $C = 29$).

\textbf{H3 (Optimal config depends on graph size) - Confirmed.} Window sizes, refresh
strategies, and start bins vary with graph size, confirming that time slicing must be tuned
per problem scale.

\textbf{H4 (Start time impacts quality) - Confirmed with caution.} Start bin 3 performs
best, though this may be influenced by experimental setup. To confirm the truly optimal
configuration, a full grid search over all parameters would be required. Nevertheless, the
results clearly indicate that the start time bin has a noticeable impact on performance.

\textbf{H5 (Time slicing reduces cost) - Rejected.} Time slicing adds overhead but
improves solution quality.

As part of the experimentation process, the current best decoder enhancement (Temp-MLP) was
combined with the best-performing time slicing configuration to investigate whether the
combination leads to further improvements. The best time slicing configuration at $C = 49$
($W = 3$, start bin $= 3$, periodic refresh interval $= 0.25$) was combined with Temp-MLP
and trained for 100 epochs across three seeds. Table~\ref{tab:temp_mlp_combo} presents the
results.

\begin{table}[H]
\centering
\caption{Temp-MLP and Time Slicing combination results ($C = 49$, 100 epochs, batch size 256).}
\label{tab:temp_mlp_combo}
\begin{adjustbox}{max width=\textwidth}
\begin{tabular}{lll}
\toprule
\textbf{Configuration} & \textbf{Mean $\pm$ SD (min)} & \textbf{$\Delta$ vs Baseline} \\
\midrule
Baseline                                    & 1095.93 $\pm$ 73.34 & -- \\
Temp-MLP standalone                         & 1041.57 $\pm$ 19.61 & $-4.96\%$ \\
Time Slicing only ($W=3$, start=3, Periodic 0.25) & 1082.81 $\pm$ 26.19 & $-1.19\%$ \\
Temp-MLP + Time Slicing                     & 1063.75 $\pm$ 6.47  & $-2.94\%$ \\
\bottomrule
\end{tabular}
\end{adjustbox}
\end{table}

The combination achieves a 2.94\% improvement over the baseline, confirming that both
extensions individually contribute meaningful gains. However, the combined configuration does
not outperform Temp-MLP standalone (1063.75 min vs 1041.57 min). This is consistent with the
pattern observed across other combinations throughout this chapter, where individual components
consistently outperform their combinations. Notably, the combination exhibits the lowest
variance of any configuration tested ($\pm$6.47 min). This suggests that combining the
encoder-level time slicing modification with the decoder-level Temp-MLP produces more stable,
but not necessarily better, solutions. The increased stability may be attributed to the
complementary nature of the two modifications. Time slicing reduces encoder noise, while
Temp-MLP regularises the decoder's exploration. Together, they constrain the solution space
more tightly across different random initialisations.

In summary, time slicing is most valuable as a solution quality enhancer at larger problem
sizes, with gains of up to 1.19\% at $C = 49$. The forward window ($W = 9$ from start
bin 3) is the most broadly effective configuration. Periodic refresh is generally preferable,
though one-time refresh remains a competitive and practical alternative.


\section{Graph-Size-Independent Architecture}
\label{sec:graph_size_indep}

The baseline M1 model has a significant practical limitation in that it was originally
designed for fixed customer set sizes. To address this, an adaptive average pooling layer is
introduced before each projection, mapping any input of arbitrary length to a fixed
128-dimensional representation (Section~\ref{sec:graph_size_independent}).

This section evaluates the pooling fix across two dimensions: (1) within-size performance,
and (2) cross-size generalisation.

\textbf{Within-Size Performance.} Table~\ref{tab:pooling_combined} compares the baseline and pooling models at $C = 19$ and
$C = 29$.

\begin{table}[H]
\centering
\caption{Within-size and cross-size performance (100 epochs, batch size 256).}
\label{tab:pooling_combined}
\begin{adjustbox}{max width=\textwidth}
\begin{tabular}{llllll}
\toprule
\textbf{Model} & \textbf{Train} & \textbf{Test}
    & \textbf{Mean $\pm$ SD (min)} & \textbf{$\Delta$ vs.\ Baseline} \\
\midrule
\multicolumn{4}{l}{\textit{Within-size}} & \\
Baseline & 19 & 19 & 565.50 $\pm$ 0.58 & -- \\
Pooling  & 19 & 19 & 570.51 $\pm$ 2.99 & $+0.89\%$ \\
Baseline & 29 & 29 & 766.79 $\pm$ 7.39 & -- \\
Pooling  & 29 & 29 & 768.12 $\pm$ 7.87 & $+0.17\%$ \\
\multicolumn{4}{l}{\textit{Cross-size}} & \\
Pooling  & 19 & 29 & 770.15 $\pm$ 10.54 & $+0.44\%$ \\
Pooling  & 29 & 19 & 572.52 $\pm$ 1.57  & $+1.24\%$ \\
\bottomrule
\end{tabular}
\end{adjustbox}
\end{table}

At $C = 19$, the pooling model underperforms by 5.01 minutes (0.89\%). At $C = 29$, the
difference is only 1.33 minutes (0.17\%). Overall, the adaptive pooling modification
introduces only a minimal performance degradation of less than 0.89\%.

\textbf{Cross-Size Generalisation.} The key innovation of the pooling is that the model can be trained on one graph size and
evaluated directly on a different graph size without retraining. Results are included in
Table~\ref{tab:pooling_combined} above.

Training on $C = 19$ and evaluating on $C = 29$ yields 770.15 minutes (0.44\% degradation
vs.\ baseline). Training on $C = 29$ and evaluating on $C = 19$ yields 572.52 minutes
(1.24\% degradation).

\textbf{Summary.} The adaptive pooling modification enables graph-size independence. A single trained instance
of the modified M1 model can be applied to any graph size without retraining, with a
solution quality degradation of less than 1\%.

\section{Beam Search}
\label{sec:beam_search}

All experiments in previous sections use greedy decoding, in which the decoder selects the
highest-probability node at each step. Beam search addresses this limitation by maintaining
$B$ candidate partial tours at each step, retaining only the $B$ most promising sequences.
At the end of decoding, the best complete tour among the $B$ candidates is returned.

Beam width $B = 3$ is evaluated (Table~\ref{tab:beam_search}).

\begin{table}[H]
\centering
\caption{Beam search results -- solution quality ($B = 3$).}
\label{tab:beam_search}
\begin{adjustbox}{max width=\textwidth}
\begin{tabular}{llll}
\toprule
\textbf{Decoding} & \textbf{Graph Size} & \textbf{Mean $\pm$ SD (min)}
    & \textbf{$\Delta$ vs.\ Greedy} \\
\midrule
Greedy     & $C = 19$ & 574.73 $\pm$ 10.78  & --- \\
Beam $B=3$ & $C = 19$ & 569.59 $\pm$ 1.03   & $-0.89\%$ \\
Greedy     & $C = 49$ & 1095.93 $\pm$ 73.34 & --- \\
Beam $B=3$ & $C = 49$ & 1095.49 $\pm$ 73.11 & $-0.04\%$ \\
\bottomrule
\end{tabular}
\end{adjustbox}
\end{table}

At $C = 19$, beam search improves performance by 5.14 minutes (0.89\%) over greedy decoding.
The cross-seed variance also drops sharply (SD = 1.03 vs.\ 10.78). This suggests that, at
this scale, beam search (with width $B = 3$) yields a meaningful improvement in solution
quality.

At $C = 49$, the improvement is negligible, at only 0.44 minutes (0.04\%). This suggests
that the search space at this scale is too large for a narrow beam to explore meaningfully.
A computational efficiency comparison between greedy and beam search decoding is left for
future work due to instrumentation constraints.

\section{Comparison with Heuristic Baselines}
\label{sec:heuristic_comparison}

This section compares the best-performing models from each above section with five
heuristics: Nearest Neighbour (NN), Greedy Edge (GE), Nearest Neighbour Random (NNR),
NN+1 lookahead, and 2-opt local search. All heuristic models are evaluated on the same
1,000 test instances used in the training process (each epoch's validation process at the
end). The same traffic data with the same cubic spline interpolation method is also used.
All results are reported as mean $\pm$ standard deviation across three random seeds (1234, 5678, 9012) for neural models, ensuring robust and reliable comparisons.

\textbf{$C = 19$.} Table~\ref{tab:heuristic_c19} presents the results at $C = 19$.

\begin{table}[H]
\centering
\caption{Solution quality and inference time comparison at $C = 19$ (mean $\pm$ SD across three seeds).}
\label{tab:heuristic_c19}
\begin{adjustbox}{max width=\textwidth}
\begin{tabular}{lll}
\toprule
\textbf{Model / Heuristic} & \textbf{Mean $\pm$ SD (min)} & \textbf{Time per Instance} \\
\midrule
\multicolumn{3}{l}{\textit{Neural Models}} \\
Baseline (M1)                 & 574.73 $\pm$ 10.78 & 0.24\,ms \\
Temp-MLP                      & 568.68 $\pm$ 2.12  & 0.36\,ms \\
Time Slicing ($W=9$, start=3) & 572.55 $\pm$ 2.86  & 0.25\,ms \\
\midrule
\multicolumn{3}{l}{\textit{Classical Heuristics}} \\
2-opt                  & 575.65 $\pm$ 89.46  & 7{,}351.47\,ms \\
Nearest Neighbour (NN) & 630.54 $\pm$ 105.49 & 9.80\,ms \\
Greedy Edge            & 633.71 $\pm$ 105.58 & 10.58\,ms \\
NNR                    & 636.38 $\pm$ 101.28 & 94.58\,ms \\
NN+1                   & 648.86 $\pm$ 104.86 & 110.21\,ms \\
\bottomrule
\end{tabular}
\end{adjustbox}
\smallskip
\begin{minipage}{\linewidth}
\small Neural model results are averaged across three random seeds (1234, 5678, 9012). Heuristic baselines are deterministic and evaluated once on the same fixed test set. Cost-Aware Gating is excluded from this table as it did not outperform the baseline at $C = 19$ (see Section~\ref{sec:cost_aware_exp} for full results).
\end{minipage}
\end{table}


\medskip

All neural models surpass the classical heuristics at $C = 19$. The baseline M1 outperforms
2-opt by 1 minute while running approximately 30,000 times faster (0.24\,ms vs 7,351.47\,ms).
Temp-MLP achieves the best neural performance (568.68 $\pm$ 2.12 min), outperforming 2-opt by
approximately 7 minutes with similarly fast inference. The remaining heuristics (NN, GE, NNR,
NN+1) all produce performance more than 55 minutes worse compared to the neural models.
Cost-Aware Gating is excluded from this table as it did not outperform the baseline at $C = 19$
(see Section~\ref{sec:cost_aware_exp} for full results).

\textbf{$C = 49$.} Table~\ref{tab:heuristic_c49} presents the results at $C = 49$.

\begin{table}[H]
\centering
\caption{Solution quality and inference time comparison at $C = 49$ (mean $\pm$ SD across three seeds).}
\label{tab:heuristic_c49}
\begin{adjustbox}{max width=\textwidth}
\begin{tabular}{lll}
\toprule
\textbf{Model / Heuristic} & \textbf{Mean $\pm$ SD (min)} & \textbf{Time per Instance} \\
\midrule
\multicolumn{3}{l}{\textit{Neural Models}} \\
Baseline (M1)                 & 1095.93 $\pm$ 73.34 & 0.64\,ms \\
Temp-MLP                      & 1041.57 $\pm$ 19.61 & 0.94\,ms \\
Cost-Aware (NN, $\lambda=0.25$) & 1055.75 $\pm$ 25.14 & 6.05\,ms \\
Time Slicing ($W=3$, start=3) & 1082.81 $\pm$ 26.19 & 0.65\,ms \\
\midrule
\multicolumn{3}{l}{\textit{Classical Heuristics}} \\
2-opt                  & 1117.61 $\pm$ 91.79  & 91{,}349.73\,ms \\
Greedy Edge            & 1154.21 $\pm$ 86.93  & 27.59\,ms \\
Nearest Neighbour (NN) & 1168.03 $\pm$ 100.46 & 23.83\,ms \\
NN+1                   & 1201.84 $\pm$ 96.65  & 630.93\,ms \\
NNR                    & 1211.52 $\pm$ 93.24  & 250.65\,ms \\
\bottomrule
\end{tabular}
\end{adjustbox}
\smallskip
\begin{minipage}{\linewidth}
\small Neural model results are averaged across three random seeds (1234, 5678, 9012). Heuristic baselines are deterministic and evaluated once on the same fixed test set.
\end{minipage}
\end{table}


\medskip

At $C = 49$, the neural models show a clear advantage over all heuristic models. Temp-MLP
achieves the best performance (1041.57 $\pm$ 19.61 min), outperforming 2-opt by approximately
76 minutes (6.8\%) while running around 97,000 times faster (0.94\,ms vs 91,349\,ms). The
baseline M1 (1095.93 $\pm$ 73.34 min) also outperforms all heuristics, beating 2-opt by
approximately 22 minutes. All proposed extensions outperform every classical heuristic at
this graph size.

\textbf{Summary.} For both graph sizes ($C = 19$ and $C = 49$), the neural models consistently outperform
classical heuristics in solution quality while offering significantly faster inference. The
advantage grows with problem size, with neural models outperforming 2-opt by 1.21\% at
$C = 19$ and up to 6.8\% at $C = 49$. These results confirm that the proposed deep
reinforcement learning approach achieves both accuracy and scalability advantages over
classical methods.

\section{Summary, Discussion, and Recommendation}
\label{sec:exp_summary_final}

This chapter evaluated four different architectural modifications to the M1 baseline on the
deterministic TDTSP with a fixed customer set. The chapter includes the key findings,
discussion points, and practical recommendations for deployment.

\subsection{Summary of Results}
\label{subsec:summary_of_results}

\textbf{Step-MLP and Decoder Variants.}
Among all decoder modifications, the Temperature MLP (Temp-MLP) is the standout performer
for both graph sizes. At $C = 19$, it reduces the mean cost by 6 minutes, and at $C = 49$,
it achieves a 4.96\% improvement. In contrast, the post-attention glimpse refinement (Option
A) degrades performance at $C = 49$, suggesting that it is better to retain the original
linear projection rather than applying a nonlinear transformation to the glimpse vector. The
Step-MLP and Option C produce better results at $C = 49$, but combining them with the
Temp-MLP reduces performance. This implies that individual components show better performance
compared to combinations.

\textbf{Cost-Aware Gating.}
This modification is clearly dependent on the graph size. At $C = 19$, the cost-aware gating
mechanism did not outperform the baseline, as none of the configurations tested were able to
beat it. However, at $C = 49$, standalone cost-aware gating with fixed $\lambda = 0.25$ and
the linear NN heuristic produced a meaningful 3.67\% improvement. An additional key finding
at $C = 19$ was that fixed lambda outperforms learnable lambda. Combining cost-aware gating
with other components degraded performance in all tested combinations.

\textbf{Time Slicing with Safe Refresh.}
Time slicing consistently adds training overhead (approximately 2.5\% at $C = 19$ and 10.7\%
at $C = 49$), but improves solution quality at both graph sizes (0.38\% and 1.19\%,
respectively). The optimal time-slicing configurations depend on the graph size. The best
configuration is $W = 3$ with one-time refresh at $C = 20$, and $W = 9$ with periodic
refresh (interval = 0.25) at $C = 19$ and $C = 49$. In most cases, periodic refresh of
the encoder improves results.

\textbf{Graph-Size-Independent Architecture.}
Introducing the adaptive pooling modification successfully mitigates the cross-size
generalisation limitation with minimal quality degradation (at most 1.24\% relative to
size-specific baselines).

\textbf{Comparison with Heuristics.}
All best-performing neural configurations outperform all classical heuristics in mean solution
quality at both $C = 19$ and $C = 49$. The 2-opt model is the most competitive among the
heuristic models, but requires a substantially longer inference time compared to the neural
models.

\subsection{Overall Recommendation}
\label{subsec:overall_recommendation}

Based on the experimental findings, the following practical recommendations are made:

\begin{itemize}
    \item \textbf{For small problem sizes ($C = 19/20$):} The Temp-MLP standalone is the
    recommended extension. Time slicing with $W = 9$, start = 3, and periodic refresh
    (interval = 0.25) also shows better performance, but the Temp-MLP standalone is highly
    recommended at this scale. Cost-aware gating should be avoided at this scale.

    \item \textbf{For larger problem sizes ($C = 49$):} The Temp-MLP remains the best
    single-component choice. However, standalone cost-aware gating (NN, fixed $\lambda =
    0.25$) and time slicing with periodic refresh ($W = 3$, interval = 0.25) also show
    better performance compared to the baseline and are worth considering.

    \item \textbf{Adaptive pooling} mitigates the graph-size dependency with a quality
    trade-off of at most 1.24\% compared to the size-specific baseline, and is worth
    considering in future developments.

    \item \textbf{Individual components} may show better performance, but combining them
    does not guarantee further gains. Therefore, it is advisable to evaluate components
    individually before finalising a configuration.
\end{itemize}


